\noindent\textbf{1.} Los tipos de ángulos que son opuestos por el vértice son:
\begin{itemize}
    \item[a)] Iguales
    \item[b)] Complementarios
    \item[c)] Suplementarios
    \item[d)] Obtusos
\end{itemize}

\noindent\textbf{2.} Dos rectas paralelas son cortadas por una transversal. Si uno de los ángulos internos mide $42^\circ$, ¿cuánto mide el ángulo alterno interno a este ángulo?
\begin{itemize}
    \item[a)] $38^\circ$
    \item[b)] $42^\circ$
    \item[c)] $128^\circ$
    \item[d)] $132^\circ$
\end{itemize}

\noindent\textbf{3.} Dos rectas paralelas son cortadas por una transversal. Si un ángulo externo $A$ mide el doble que su ángulo colateral externo $B$, ¿cuánto mide el ángulo $B$?
\begin{itemize}
    \item[a)] $30^\circ$
    \item[b)] $60^\circ$
    \item[c)] $90^\circ$
    \item[d)] $120^\circ$
\end{itemize}

\noindent\textbf{4.} Dos rectas paralelas son cortadas por una transversal. ¿Cómo se llama la pareja de ángulos que están del mismo lado de la transversal, pero uno está en el interior de las paralelas, el otro está por fuera y no son adyacentes?
\begin{itemize}
    \item[a)] Colaterales externos
    \item[b)] Opuestos externos
    \item[c)] Alternos internos
    \item[d)] Correspondientes
\end{itemize}

\noindent\textbf{5.} ¿Cuál de las siguientes parejas de ángulos suman $180^\circ$?
\begin{itemize}
    \item[a)] Correspondientes
    \item[b)] Opuestos por el vértice
    \item[c)] Colaterales internos
    \item[d)] Alternos externos
\end{itemize}

\noindent\textbf{6.} ¿Cuántos pares de lados paralelos tiene un trapecio escaleno?
\begin{itemize}
    \item[a)] Tiene 2
    \item[b)] Tiene 1
    \item[c)] No tiene parejas de lados paralelos
    \item[d)] Depende de sus medidas
\end{itemize}

\noindent\textbf{7.} ¿Cuántas parejas de lados paralelos tiene un romboide?

\begin{itemize}
    \item[a)] Tiene 2
    \item[b)] Tiene 1
    \item[c)] No tiene parejas de lados paralelos
    \item[d)] Depende de sus medidas
\end{itemize}


\noindent\textbf{8.} ¿Cuántas parejas de lados paralelos tiene un trapezoide?

\begin{itemize}
    \item[a)] Tiene 2
    \item[b)] Tiene 1
    \item[c)] No tiene parejas de lados paralelos
    \item[d)] Depende de sus medidas
\end{itemize}

\noindent\textbf{9.} En un rectángulo la base mide el triple de su altura. Si su perímetro es de 80 metros, ¿cuánto mide su área?

\begin{itemize}
    \item[a)] $100\,\text{m}^2$
    \item[b)] $200\,\text{m}^2$
    \item[c)] $240\,\text{m}^2$
    \item[d)] $300\,\text{m}^2$
\end{itemize}

\noindent\textbf{8.} ¿Cuántas parejas de lados paralelos tiene un trapezoide?

\begin{itemize}
    \item[a)] Tiene 2
    \item[b)] Tiene 1
    \item[c)] No tiene parejas de lados paralelos
    \item[d)] Depende de sus medidas
\end{itemize}

\noindent\textbf{9.} En un rectángulo la base mide el triple de su altura. Si su perímetro es de 80 metros, ¿cuánto mide su área?

\begin{itemize}
    \item[a)] $100\,\text{m}^2$
    \item[b)] $200\,\text{m}^2$
    \item[c)] $240\,\text{m}^2$
    \item[d)] $300\,\text{m}^2$
\end{itemize}

\noindent\textbf{10.} En un trapecio rectangular la altura mide 6 metros, la base mayor mide 2 metros más que la base menor y su área es de 36 metros cuadrados. ¿Cuánto mide la base menor?

\begin{itemize}
    \item[a)] 5 metros
    \item[b)] 2 metros
    \item[c)] 10 metros
    \item[d)] 12 metros
\end{itemize}

\noindent\textbf{11.} ¿Cómo se llama el polígono regular que tiene 9 lados?

\begin{itemize}
    \item[a)] Eneágono
    \item[b)] Decágono
    \item[c)] Heptágono
    \item[d)] Endecágono
\end{itemize}

\noindent\textbf{12.} La suma de los ángulos interiores de un polígono es de $720^\circ$. ¿Cómo se llama ese polígono?

\begin{itemize}
    \item[a)] Hexágono
    \item[b)] Pentágono
    \item[c)] Heptágono
    \item[d)] Eneágono
\end{itemize}

\noindent\textbf{13.} ¿Cuánto miden los ángulos centrales de un polígono regular de 3 lados?

\begin{itemize}
    \item[a)] $60^\circ$
    \item[b)] $90^\circ$
    \item[c)] $120^\circ$
    \item[d)] $180^\circ$
\end{itemize}

\noindent\textbf{14.} La medida de un ángulo interno de un polígono regular es de $120^\circ$. ¿Cuántos lados tiene?

\begin{itemize}
    \item[a)] 6
    \item[b)] 8
    \item[c)] 10
    \item[d)] 12
\end{itemize}

\noindent\textbf{15.} La longitud de cada lado de un pentágono regular es de 2 metros. Si el área del pentágono es de 25 metros cuadrados, basándose en la fórmula de área para polígonos regulares, ¿cuánto mide su apotema?

\begin{itemize}
    \item[a)] 2.5 metros
    \item[b)] 4 metros
    \item[c)] 5 metros
    \item[d)] 6.5 metros
\end{itemize}

\noindent\textbf{16.} ¿Cómo se llama el segmento que une dos puntos de una circunferencia, pero no pasa por su centro?

\begin{itemize}
    \item[a)] Arco
    \item[b)] Secante
    \item[c)] Radio
    \item[d)] Cuerda
\end{itemize}


\newpage

\section*{Tabla de respuestas}

\begin{center}
\begin{tabular}{|c|c|l|}
\hline
\textbf{Pregunta} & \textbf{Opción} & \textbf{Respuesta} \\ \hline
1  & a) & Iguales \\ \hline
2  & b) & $42^\circ$ \\ \hline
3  & b) & $60^\circ$ \\ \hline
4  & d) & Correspondientes \\ \hline
5  & c) & Colaterales internos \\ \hline
6  & b) & Tiene 1 \\ \hline
7  & a) & Tiene 2 \\ \hline
8  & c) & Ningún par de lados paralelos \\ \hline
9  & d) & $300\,\text{m}^2$ \\ \hline
10 & a) & 5 metros  \\ \hline
11 & a) & Eneágono \\ \hline
12 & a) & Hexágono \\ \hline
13 & c) & $120^\circ$ \\ \hline
14 & a) & 6 lados \\ \hline
15 & c) & 5 metros \\ \hline
16 & d) & Cuerda \\ \hline
\end{tabular}
\end{center}


\newpage

\section*{Examen de Geometría -- Versión B}

\noindent\textbf{1.} Dos ángulos que suman $90^\circ$ reciben el nombre de:

\begin{itemize}
    \item[a)] Suplementarios
    \item[b)] Complementarios
    \item[c)] Opuestos por el vértice
    \item[d)] Obtusos
\end{itemize}


\noindent\textbf{2.} Dos rectas paralelas son cortadas por una transversal. Si uno de los ángulos internos mide $65^\circ$, ¿cuánto mide el ángulo alterno interno a este ángulo?

\begin{itemize}
    \item[a)] $115^\circ$
    \item[b)] $25^\circ$
    \item[c)] $65^\circ$
    \item[d)] $125^\circ$
\end{itemize}


\noindent\textbf{3.} Dos rectas paralelas son cortadas por una transversal. Si un ángulo externo $A$ mide el triple que su ángulo colateral externo $B$, ¿cuánto mide el ángulo $B$?

\begin{itemize}
    \item[a)] $45^\circ$
    \item[b)] $60^\circ$
    \item[c)] $90^\circ$
    \item[d)] $135^\circ$
\end{itemize}


\noindent\textbf{4.} Dos rectas paralelas son cortadas por una transversal. ¿Cómo se llama la pareja de ángulos que están entre las rectas paralelas, en lados opuestos de la transversal y no son adyacentes?

\begin{itemize}
    \item[a)] Correspondientes
    \item[b)] Colaterales externos
    \item[c)] Alternos internos
    \item[d)] Opuestos por el vértice
\end{itemize}


\noindent\textbf{5.} ¿Cuál de las siguientes parejas de ángulos tienen la misma medida cuando dos rectas paralelas son cortadas por una transversal?

\begin{itemize}
    \item[a)] Colaterales internos
    \item[b)] Colaterales externos
    \item[c)] Suplementarios
    \item[d)] Alternos externos
\end{itemize}


\noindent\textbf{6.} ¿Cuántos pares de lados paralelos tiene un trapecio isósceles?

\begin{itemize}
    \item[a)] Tiene 2
    \item[b)] Tiene 1
    \item[c)] No tiene parejas de lados paralelos
    \item[d)] Depende de sus medidas
\end{itemize}


\noindent\textbf{7.} ¿Cuántas parejas de lados paralelos tiene un rombo?

\begin{itemize}
    \item[a)] Tiene 1
    \item[b)] No tiene parejas de lados paralelos
    \item[c)] Tiene 2
    \item[d)] Depende de sus medidas
\end{itemize}


\noindent\textbf{8.} ¿Cuántas parejas de lados paralelos tiene un trapezoide irregular?

\begin{itemize}
    \item[a)] Tiene 1
    \item[b)] Tiene 2
    \item[c)] Depende de sus medidas
    \item[d)] No tiene parejas de lados paralelos
\end{itemize}


\noindent\textbf{9.} En un rectángulo la base mide el doble de su altura. Si su perímetro es de 60 metros, ¿cuánto mide su área?

\begin{itemize}
    \item[a)] $100\,\text{m}^2$
    \item[b)] $150\,\text{m}^2$
    \item[c)] $200\,\text{m}^2$
    \item[d)] $300\,\text{m}^2$
\end{itemize}


\noindent\textbf{10.} En un trapecio rectangular la altura mide 4 metros, la base mayor mide 4 metros más que la base menor y su área es de 32 metros cuadrados. ¿Cuánto mide la base menor?

\begin{itemize}
    \item[a)] 4 metros
    \item[b)] 6 metros
    \item[c)] 8 metros
    \item[d)] 12 metros
\end{itemize}


\noindent\textbf{11.} ¿Cómo se llama el polígono regular que tiene 8 lados?

\begin{itemize}
    \item[a)] Hexágono
    \item[b)] Octágono
    \item[c)] Eneágono
    \item[d)] Decágono
\end{itemize}


\noindent\textbf{12.} La suma de los ángulos interiores de un polígono es de $900^\circ$. ¿Cómo se llama ese polígono?

\begin{itemize}
    \item[a)] Pentágono
    \item[b)] Hexágono
    \item[c)] Heptágono
    \item[d)] Octágono
\end{itemize}


\noindent\textbf{13.} ¿Cuánto miden los ángulos centrales de un polígono regular de 6 lados?

\begin{itemize}
    \item[a)] $45^\circ$
    \item[b)] $60^\circ$
    \item[c)] $90^\circ$
    \item[d)] $120^\circ$
\end{itemize}


\noindent\textbf{14.} La medida de un ángulo interno de un polígono regular es de $135^\circ$. ¿Cuántos lados tiene?

\begin{itemize}
    \item[a)] 5
    \item[b)] 6
    \item[c)] 8
    \item[d)] 10
\end{itemize}


\noindent\textbf{15.} La longitud de cada lado de un hexágono regular es de 4 metros. Si el área del hexágono es de 60 metros cuadrados, basándose en la fórmula de área para polígonos regulares, ¿cuánto mide su apotema?

\begin{itemize}
    \item[a)] 3 metros
    \item[b)] 4 metros
    \item[c)] 5 metros
    \item[d)] 6 metros
\end{itemize}


\noindent\textbf{16.} ¿Cómo se llama el segmento que une el centro de una circunferencia con cualquier punto de ella?

\begin{itemize}
    \item[a)] Diámetro
    \item[b)] Cuerda
    \item[c)] Secante
    \item[d)] Radio
\end{itemize}


\newpage

\section*{Tabla de respuestas -- Versión B}

\begin{center}
\begin{tabular}{|c|c|l|}
\hline
\textbf{Pregunta} & \textbf{Opción} & \textbf{Respuesta} \\ \hline
1  & b) & Complementarios \\ \hline
2  & c) & $65^\circ$ \\ \hline
3  & a) & $45^\circ$ \\ \hline
4  & c) & Alternos internos \\ \hline
5  & d) & Alternos externos \\ \hline
6  & b) & Tiene 1 \\ \hline
7  & c) & Tiene 2 \\ \hline
8  & d) & Ningún par de lados paralelos \\ \hline
9  & c) & $200\,\text{m}^2$ \\ \hline
10 & b) & 6 metros \\ \hline
11 & b) & Octágono \\ \hline
12 & c) & Heptágono \\ \hline
13 & b) & $60^\circ$ \\ \hline
14 & c) & 8 lados \\ \hline
15 & c) & 5 metros \\ \hline
16 & d) & Radio \\ \hline
\end{tabular}
\end{center}

\newpage

\section*{Examen de Geometría -- Versión C}

\noindent\textbf{1.} ¿Cómo se llama un ángulo que mide más de $90^\circ$ pero menos de $180^\circ$?

\begin{itemize}
    \item[a)] Ángulo agudo
    \item[b)] Ángulo recto
    \item[c)] Ángulo obtuso
    \item[d)] Ángulo llano
\end{itemize}


\noindent\textbf{2.} Dos ángulos son complementarios. Si uno de ellos mide $35^\circ$, ¿cuánto mide el otro?

\begin{itemize}
    \item[a)] $55^\circ$
    \item[b)] $65^\circ$
    \item[c)] $145^\circ$
    \item[d)] $155^\circ$
\end{itemize}


\noindent\textbf{3.} Dos rectas paralelas son cortadas por una transversal. Si uno de los ángulos colaterales internos mide $125^\circ$, ¿cuánto mide el otro?

\begin{itemize}
    \item[a)] $125^\circ$
    \item[b)] $65^\circ$
    \item[c)] $45^\circ$
    \item[d)] $55^\circ$
\end{itemize}


\noindent\textbf{4.} Un triángulo tiene tres lados que miden 5 cm, 5 cm y 8 cm. ¿Cómo se clasifica según sus lados?

\begin{itemize}
    \item[a)] Equilátero
    \item[b)] Isósceles
    \item[c)] Escaleno
    \item[d)] Rectángulo
\end{itemize}


\noindent\textbf{5.} En un triángulo dos de sus ángulos interiores miden $45^\circ$ y $65^\circ$. ¿Cuánto mide el tercer ángulo?

\begin{itemize}
    \item[a)] $60^\circ$
    \item[b)] $80^\circ$
    \item[c)] $70^\circ$
    \item[d)] $90^\circ$
\end{itemize}


\noindent\textbf{6.} ¿Cuál de los siguientes cuadriláteros tiene exactamente un par de lados paralelos y dos ángulos rectos?

\begin{itemize}
    \item[a)] Romboide
    \item[b)] Trapecio rectangular
    \item[c)] Rombo
    \item[d)] Trapezoide
\end{itemize}


\noindent\textbf{7.} Un cuadrado tiene lados de 7 centímetros. ¿Cuánto mide su perímetro?

\begin{itemize}
    \item[a)] 14 cm
    \item[b)] 21 cm
    \item[c)] 49 cm
    \item[d)] 28 cm
\end{itemize}


\noindent\textbf{8.} Un triángulo tiene una base de 12 metros y una altura de 5 metros. ¿Cuánto mide su área?

\begin{itemize}
    \item[a)] $30\,\text{m}^2$
    \item[b)] $60\,\text{m}^2$
    \item[c)] $17\,\text{m}^2$
    \item[d)] $24\,\text{m}^2$
\end{itemize}


\noindent\textbf{9.} Un trapecio tiene una base mayor de 10 metros, una base menor de 6 metros y una altura de 4 metros. ¿Cuánto mide su área?

\begin{itemize}
    \item[a)] $64\,\text{m}^2$
    \item[b)] $40\,\text{m}^2$
    \item[c)] $32\,\text{m}^2$
    \item[d)] $16\,\text{m}^2$
\end{itemize}


\noindent\textbf{10.} Un círculo tiene un radio de 5 centímetros. Utilizando $\pi=3.14$, ¿cuánto mide su área?

\begin{itemize}
    \item[a)] $31.4\,\text{cm}^2$
    \item[b)] $78.5\,\text{cm}^2$
    \item[c)] $15.7\,\text{cm}^2$
    \item[d)] $25\,\text{cm}^2$
\end{itemize}


\noindent\textbf{11.} ¿Cuántas diagonales se pueden trazar desde un solo vértice de un octágono?

\begin{itemize}
    \item[a)] 8
    \item[b)] 4
    \item[c)] 6
    \item[d)] 5
\end{itemize}


\noindent\textbf{12.} ¿Cuánto mide cada ángulo central de un decágono regular?

\begin{itemize}
    \item[a)] $36^\circ$
    \item[b)] $45^\circ$
    \item[c)] $60^\circ$
    \item[d)] $72^\circ$
\end{itemize}


\noindent\textbf{13.} La suma de los ángulos interiores de un polígono es de $1080^\circ$. ¿Cuántos lados tiene?

\begin{itemize}
    \item[a)] 6 lados
    \item[b)] 7 lados
    \item[c)] 8 lados
    \item[d)] 10 lados
\end{itemize}


\noindent\textbf{14.} Dos triángulos tienen sus tres lados correspondientes proporcionales. ¿Qué criterio de semejanza se utiliza?

\begin{itemize}
    \item[a)] AA (Ángulo-Ángulo)
    \item[b)] LAL (Lado-Ángulo-Lado)
    \item[c)] AAA (Ángulo-Ángulo-Ángulo)
    \item[d)] LLL (Lado-Lado-Lado)
\end{itemize}


\noindent\textbf{15.} Un pentágono regular tiene un perímetro de 30 centímetros y una apotema de 4 centímetros. ¿Cuánto mide su área?

\begin{itemize}
    \item[a)] $120\,\text{cm}^2$
    \item[b)] $60\,\text{cm}^2$
    \item[c)] $34\,\text{cm}^2$
    \item[d)] $15\,\text{cm}^2$
\end{itemize}


\noindent\textbf{16.} Un triángulo rectángulo tiene catetos de 9 centímetros y 12 centímetros. ¿Cuánto mide su hipotenusa?

\begin{itemize}
    \item[a)] 15 cm
    \item[b)] 18 cm
    \item[c)] 21 cm
    \item[d)] 25 cm
\end{itemize}


\newpage

\section*{Tabla de respuestas -- Versión C}

\begin{center}
\begin{tabular}{|c|c|l|}
\hline
\textbf{Pregunta} & \textbf{Opción} & \textbf{Respuesta} \\ \hline
1  & c) & Ángulo obtuso \\ \hline
2  & a) & $55^\circ$ \\ \hline
3  & d) & $55^\circ$ \\ \hline
4  & b) & Isósceles \\ \hline
5  & c) & $70^\circ$ \\ \hline
6  & b) & Trapecio rectangular \\ \hline
7  & d) & 28 cm \\ \hline
8  & a) & $30\,\text{m}^2$ \\ \hline
9  & c) & $32\,\text{m}^2$ \\ \hline
10 & b) & $78.5\,\text{cm}^2$ \\ \hline
11 & d) & 5 diagonales \\ \hline
12 & a) & $36^\circ$ \\ \hline
13 & c) & 8 lados \\ \hline
14 & d) & LLL \\ \hline
15 & b) & $60\,\text{cm}^2$ \\ \hline
16 & a) & 15 cm \\ \hline
\end{tabular}
\end{center}

\newpage

\section*{Examen de Geometría -- Versión D}

\noindent\textbf{1.} ¿Cómo se clasifica un ángulo que mide $235^\circ$?

\begin{itemize}
    \item[a)] Ángulo agudo
    \item[b)] Ángulo cóncavo o reflejo
    \item[c)] Ángulo obtuso
    \item[d)] Ángulo completo
\end{itemize}


\noindent\textbf{2.} Dos ángulos son complementarios. Si uno mide $27^\circ$, ¿cuánto mide el otro?

\begin{itemize}
    \item[a)] $63^\circ$
    \item[b)] $153^\circ$
    \item[c)] $27^\circ$
    \item[d)] $117^\circ$
\end{itemize}


\noindent\textbf{3.} Dos ángulos son conjugados. Si uno mide $135^\circ$, ¿cuánto mide el otro?

\begin{itemize}
    \item[a)] $45^\circ$
    \item[b)] $135^\circ$
    \item[c)] $315^\circ$
    \item[d)] $225^\circ$
\end{itemize}


\noindent\textbf{4.} Dos rectas se cruzan. Si uno de los ángulos mide $74^\circ$, ¿cuánto mide su opuesto por el vértice?

\begin{itemize}
    \item[a)] $106^\circ$
    \item[b)] $286^\circ$
    \item[c)] $74^\circ$
    \item[d)] $148^\circ$
\end{itemize}


\noindent\textbf{5.} Dos rectas paralelas son cortadas por una transversal. Si un ángulo mide $118^\circ$, ¿cuánto mide su correspondiente?

\begin{itemize}
    \item[a)] $118^\circ$
    \item[b)] $62^\circ$
    \item[c)] $124^\circ$
    \item[d)] $242^\circ$
\end{itemize}


\noindent\textbf{6.} Dos rectas paralelas son cortadas por una transversal. Si un ángulo colateral externo mide $146^\circ$, ¿cuánto mide el otro?

\begin{itemize}
    \item[a)] $146^\circ$
    \item[b)] $44^\circ$
    \item[c)] $214^\circ$
    \item[d)] $34^\circ$
\end{itemize}


\noindent\textbf{7.} Un triángulo tiene lados de 5 cm, 8 cm y 10 cm. ¿Cómo se clasifica según sus lados?

\begin{itemize}
    \item[a)] Isósceles
    \item[b)] Equilátero
    \item[c)] Escaleno
    \item[d)] Rectángulo
\end{itemize}


\noindent\textbf{8.} Un triángulo tiene ángulos interiores de $100^\circ$, $42^\circ$ y $38^\circ$. ¿Cómo se clasifica según sus ángulos?

\begin{itemize}
    \item[a)] Acutángulo
    \item[b)] Obtusángulo
    \item[c)] Rectángulo
    \item[d)] Equilátero
\end{itemize}


\noindent\textbf{9.} En un triángulo, dos ángulos interiores miden $47^\circ$ y $68^\circ$. ¿Cuánto mide el tercer ángulo?

\begin{itemize}
    \item[a)] $65^\circ$
    \item[b)] $115^\circ$
    \item[c)] $75^\circ$
    \item[d)] $55^\circ$
\end{itemize}


\noindent\textbf{10.} En un triángulo, uno de sus ángulos interiores mide $55^\circ$. ¿Cuánto mide el ángulo exterior adyacente?

\begin{itemize}
    \item[a)] $55^\circ$
    \item[b)] $109^\circ$
    \item[c)] $125^\circ$
    \item[d)] $180^\circ$
\end{itemize}


\noindent\textbf{11.} ¿Qué cuadrilátero tiene dos pares de lados paralelos y sus cuatro lados iguales, y ninguno de sus ángulos es recto?

\begin{itemize}
    \item[a)] Cuadrado
    \item[b)] Rombo
    \item[c)] Rectángulo
    \item[d)] Romboide
\end{itemize}


\noindent\textbf{12.} ¿Cómo se llama un cuadrilátero que no tiene ningún par de lados paralelos?

\begin{itemize}
    \item[a)] Trapecio
    \item[b)] Cuadrado
    \item[c)] Romboide
    \item[d)] Trapezoide
\end{itemize}


\noindent\textbf{13.} Un trapecio tiene exactamente un par de lados paralelos y sus dos lados no paralelos son iguales. ¿De qué tipo es?

\begin{itemize}
    \item[a)] Trapecio isósceles
    \item[b)] Trapecio escaleno
    \item[c)] Trapecio rectángulo
    \item[d)] Trapezoide
\end{itemize}


\noindent\textbf{14.} ¿Cómo se llama el segmento que une dos puntos de una circunferencia sin necesidad de pasar por su centro?

\begin{itemize}
    \item[a)] Radio
    \item[b)] Diámetro
    \item[c)] Cuerda
    \item[d)] Arco
\end{itemize}


\noindent\textbf{15.} ¿Cómo se llama la recta que toca una circunferencia en un solo punto?

\begin{itemize}
    \item[a)] Recta secante
    \item[b)] Recta tangente
    \item[c)] Diámetro
    \item[d)] Radio
\end{itemize}


\noindent\textbf{16.} Un rectángulo tiene 9 cm de base y 4 cm de altura. ¿Cuál es su perímetro?

\begin{itemize}
    \item[a)] 36 cm
    \item[b)] 13 cm
    \item[c)] 52 cm
    \item[d)] 26 cm
\end{itemize}


\noindent\textbf{17.} Un cuadrado tiene lados de 6 cm. ¿Cuál es su área?

\begin{itemize}
    \item[a)] $36\,\text{cm}^2$
    \item[b)] $24\,\text{cm}^2$
    \item[c)] $12\,\text{cm}^2$
    \item[d)] $18\,\text{cm}^2$
\end{itemize}


\noindent\textbf{18.} Un triángulo tiene base de 14 m y altura de 9 m. ¿Cuál es su área?

\begin{itemize}
    \item[a)] $126\,\text{m}^2$
    \item[b)] $46\,\text{m}^2$
    \item[c)] $63\,\text{m}^2$
    \item[d)] $23\,\text{m}^2$
\end{itemize}


\noindent\textbf{19.} Un trapecio tiene base mayor de 14 cm, base menor de 8 cm y altura de 5 cm. ¿Cuál es su área?

\begin{itemize}
    \item[a)] $110\,\text{cm}^2$
    \item[b)] $55\,\text{cm}^2$
    \item[c)] $22\,\text{cm}^2$
    \item[d)] $35\,\text{cm}^2$
\end{itemize}


\noindent\textbf{20.} Un círculo tiene radio de 4 cm. Utilizando $\pi=3.14$, ¿cuál es su área?

\begin{itemize}
    \item[a)] $12.56\,\text{cm}^2$
    \item[b)] $25.12\,\text{cm}^2$
    \item[c)] $16\,\text{cm}^2$
    \item[d)] $50.24\,\text{cm}^2$
\end{itemize}


\noindent\textbf{21.} Una circunferencia tiene radio de 7 cm. Utilizando $\pi=3.14$, ¿cuál es su longitud?

\begin{itemize}
    \item[a)] 43.96 cm
    \item[b)] 153.86 cm
    \item[c)] 21.98 cm
    \item[d)] 14 cm
\end{itemize}


\noindent\textbf{22.} En un polígono regular, ¿qué es la apotema?

\begin{itemize}
    \item[a)] El segmento que une el centro con un vértice
    \item[b)] El segmento que une dos vértices no consecutivos
    \item[c)] El segmento perpendicular que une el centro con el punto medio de un lado
    \item[d)] El segmento que une dos vértices consecutivos
\end{itemize}


\noindent\textbf{23.} Un heptágono regular tiene lados de 5 cm. ¿Cuánto mide su perímetro?

\begin{itemize}
    \item[a)] 12 cm
    \item[b)] 35 cm
    \item[c)] 25 cm
    \item[d)] 40 cm
\end{itemize}


\noindent\textbf{24.} Un polígono regular tiene perímetro de 48 cm y apotema de 3.5 cm. ¿Cuál es su área?

\begin{itemize}
    \item[a)] $168\,\text{cm}^2$
    \item[b)] $51.5\,\text{cm}^2$
    \item[c)] $24\,\text{cm}^2$
    \item[d)] $84\,\text{cm}^2$
\end{itemize}


\noindent\textbf{25.} ¿Cuántas diagonales se pueden trazar desde un solo vértice de un decágono?

\begin{itemize}
    \item[a)] 7
    \item[b)] 10
    \item[c)] 14
    \item[d)] 35
\end{itemize}


\noindent\textbf{26.} ¿Cuántas diagonales tiene en total un octágono?

\begin{itemize}
    \item[a)] 8
    \item[b)] 5
    \item[c)] 20
    \item[d)] 40
\end{itemize}


\noindent\textbf{27.} ¿Cuánto mide cada ángulo central de un dodecágono regular?

\begin{itemize}
    \item[a)] $12^\circ$
    \item[b)] $30^\circ$
    \item[c)] $36^\circ$
    \item[d)] $45^\circ$
\end{itemize}


\noindent\textbf{28.} ¿Cuánto suman los ángulos interiores de un eneágono de 9 lados?

\begin{itemize}
    \item[a)] $1440^\circ$
    \item[b)] $1080^\circ$
    \item[c)] $1620^\circ$
    \item[d)] $1260^\circ$
\end{itemize}


\noindent\textbf{29.} Dos triángulos tienen dos ángulos correspondientes iguales. ¿Qué criterio permite afirmar que son semejantes?

\begin{itemize}
    \item[a)] AA (Ángulo-Ángulo)
    \item[b)] LAL (Lado-Ángulo-Lado)
    \item[c)] LLL (Lado-Lado-Lado)
    \item[d)] Ninguno de los anteriores
\end{itemize}


\noindent\textbf{30.} Dos triángulos tienen lados correspondientes de 4 cm y 7 cm, y de 8 cm y 14 cm; el ángulo comprendido entre esos lados mide $50^\circ$ en ambos. ¿Qué criterio de semejanza se utiliza?

\begin{itemize}
    \item[a)] AA (Ángulo-Ángulo)
    \item[b)] LAL (Lado-Ángulo-Lado)
    \item[c)] LLL (Lado-Lado-Lado)
    \item[d)] Ninguno de los anteriores
\end{itemize}


\noindent\textbf{31.} Un triángulo tiene lados de 3 cm, 5 cm y 7 cm, y otro de 6 cm, 10 cm y 14 cm. ¿Qué criterio de semejanza se cumple?

\begin{itemize}
    \item[a)] AA (Ángulo-Ángulo)
    \item[b)] LAL (Lado-Ángulo-Lado)
    \item[c)] LLL (Lado-Lado-Lado)
    \item[d)] No son semejantes
\end{itemize}


\noindent\textbf{32.} El ángulo central de un polígono regular mide $40^\circ$ y su ángulo exterior mide $(5x+10)^\circ$. ¿Cuánto vale $x$?

\begin{itemize}
    \item[a)] $x=4$
    \item[b)] $x=5$
    \item[c)] $x=8$
    \item[d)] $x=6$
\end{itemize}


\newpage

\section*{Tabla de respuestas -- Versión D}

\section*{Tabla de respuestas}

\begin{center}
\begin{tabular}{|c|c|l|}
\hline
\textbf{Pregunta} & \textbf{Opción} & \textbf{Respuesta} \\ \hline
1  & b) & Ángulo cóncavo o reflejo \\ \hline
2  & a) & $63^\circ$ \\ \hline
3  & d) & $225^\circ$ \\ \hline
4  & c) & $74^\circ$ \\ \hline
5  & a) & $118^\circ$ \\ \hline
6  & d) & $34^\circ$ \\ \hline
7  & c) & Escaleno \\ \hline
8  & b) & Obtusángulo \\ \hline
9  & a) & $65^\circ$ \\ \hline
10 & c) & $125^\circ$ \\ \hline
11 & b) & Rombo \\ \hline
12 & d) & Trapezoide \\ \hline
13 & a) & Trapecio isósceles \\ \hline
14 & c) & Cuerda \\ \hline
15 & b) & Recta tangente \\ \hline
16 & d) & 26 cm \\ \hline
17 & a) & $36\,\text{cm}^2$ \\ \hline
18 & c) & $63\,\text{m}^2$ \\ \hline
19 & b) & $55\,\text{cm}^2$ \\ \hline
20 & d) & $50.24\,\text{cm}^2$ \\ \hline
21 & a) & 43.96 cm \\ \hline
22 & c) & Apotema \\ \hline
23 & b) & 35 cm \\ \hline
24 & d) & $84\,\text{cm}^2$ \\ \hline
25 & a) & 7 diagonales \\ \hline
26 & c) & 20 diagonales \\ \hline
27 & b) & $30^\circ$ \\ \hline
28 & d) & $1260^\circ$ \\ \hline
29 & a) & AA (Ángulo-Ángulo) \\ \hline
30 & b) & LAL (Lado-Ángulo-Lado) \\ \hline
31 & c) & LLL (Lado-Lado-Lado) \\ \hline
32 & d) & $x=6$ \\ \hline
\end{tabular}
\end{center}